%
% This is the Appendix B file (appB.tex)
%
\appendix{Choosing a Muon Annulus for Mass Sensitivity}\label{app:annulus}

As presented in~\Cref{ssec:observables}, a single muon annulus was chosen to use in the mass sensitivity analysis from a total of 20 studied annuli.
The annulus spanning \qtyrange{800}{850}{\meter} from the air-shower axis was chosen by simultaneously attempting to maximize the proton-iron separation power using the FOM metric in~\Cref{eq:fom} while qualitatively minimizing the ratio of $N_{\mu}$ uncertainty-to-counts from Poisson statistics (i.e. $\sqrt{N_{\mu}}/N_{\mu}$).
The proton-iron merit factors of muon numbers within annuli of different spacings from the shower axis are shown in~\Cref{fig:muon_annuli}.
The uncertainty bands in the FOM values were estimated with the same bootstrapping method described in~\Cref{sec:fisher}.

\begin{figure}
  \centering
  \includegraphics[trim={0cm 0cm 0cm 0cm},clip,width=0.49\textwidth]{appendix/B/EPOSLHCR_MassFOM_MuonRings_zen40_50_first10rings_wDataCuts_pFe_electronAtXmaxScaling.pdf}
  \includegraphics[trim={0cm 0cm 0cm 0cm},clip,width=0.49\textwidth]{appendix/B/EPOSLHCR_MassFOM_MuonRings_zen40_50_last10rings_wDataCuts_pFe_electronAtXmaxScaling.pdf}
  \caption[Muon annulus proton-iron merit factors.]{The proton-iron figure of merit for several muon annuli ranging between \qtyrange{0}{500}{\meter} (left) and \qtyrange{500}{1000}{\meter} (right) from the air-shower axis. Each annulus has a width of \qty{50}{\meter} in the shower plane and is then projected onto the ground. Within this studied zenith range, the annuli within \qty{500}{\meter} to the shower axis quickly become less mass sensitive as the annulus distance approaches the shower axis.}
  \label{fig:muon_annuli}
\end{figure}

The muon annulus farthest from the air-shower axis (\qtyrange{950}{1000}{\meter}) exhibits the maximum proton-iron separation power amongst all energy bins.
However, there are only small differences (approximately 0.1 or less) in the proton-iron separation power between the \qtyrange{950}{1000}{\meter} annulus and other annuli far (approximately $>$\qty{800}{\meter}) from the shower axis.
The described reasoning holds for other zenith angles and hadronic interaction models.
The EPOS LHC-R proton-iron separation power for the same annuli in zenith bins of \qtyrange{0}{20}{\degree} and \qtyrange{20}{40}{\degree} is presented in~\Cref{fig:muon_annuli_zenith}.
Annuli farther from the shower axis, in general, exhibit a higher proton-iron mass separation power in most energy bins except at the highest energies, where the annuli near \qty{500}{\meter} from the shower axis provides a larger separation than the \qtyrange{950}{1000}{\meter} annulus.
The high-gain PMTs of the Auger water-Cherenkov detectors are likely to saturate at distances up to \qty{500}{\meter} from the reconstructed air-shower axis for the highest energy air showers~\cite{auger_jinst2020_sdreconstruction}, making annuli within \qty{500}{\meter} a poor practical choice for studying the mass sensitivity of a muon observable for ultra-high-energy cosmic rays.

\begin{figure}
  \centering
  \includegraphics[trim={0cm 0cm 0cm 0cm},clip,width=0.45\textwidth]{appendix/B/EPOSLHCR_MassFOM_MuonRings_zen0_20_first10rings_wDataCuts_pFe_electronAtXmaxScaling.pdf}
  \includegraphics[trim={0cm 0cm 0cm 0cm},clip,width=0.45\textwidth]{appendix/B/EPOSLHCR_MassFOM_MuonRings_zen0_20_last10rings_wDataCuts_pFe_electronAtXmaxScaling.pdf}\\
  \includegraphics[trim={0cm 0cm 0cm 0cm},clip,width=0.45\textwidth]{appendix/B/EPOSLHCR_MassFOM_MuonRings_zen20_40_first10rings_wDataCuts_pFe_electronAtXmaxScaling.pdf}
  \includegraphics[trim={0cm 0cm 0cm 0cm},clip,width=0.45\textwidth]{appendix/B/EPOSLHCR_MassFOM_MuonRings_zen20_40_last10rings_wDataCuts_pFe_electronAtXmaxScaling.pdf}
  \caption[Muon annulus merit factors for different zenith bins.]{The proton-iron figure of merit of all muon annuli, studied within zenith bins of \qtyrange{0}{20}{\degree} (top) and \qtyrange{20}{40}{\degree} (bottom). Annuli within \qtyrange{0}{500}{\meter} from the shower axis are shown on the left and those within \qtyrange{500}{1000}{\meter} are shown on the right. The proton-iron separation power of the muon annuli depends on zenith angle, with the muonic component of nearly vertical showers exhibiting higher merit factors than those of more inclined showers. However, annuli farther from the shower axis generally exhibit higher or nearly equal separation power than annuli closer to the shower axis in most energy bins.}
  \label{fig:muon_annuli_zenith}
\end{figure}

Generally, moving closer to the shower axis lessens the ratio of $N_{\mu}$ uncertainty-to-counts due to both the Poissonian statistics of counting single muons in surface particle detectors and the lateral distribution of muons at ground.
As air-shower energy decreases so does the size of the particle footprint on ground, further affecting the difference in $N_{\mu}$ uncertainty-to-counts between distances close (a few hundred meters) and far (near 1000 meters) from the air-shower axis.
Alternatively, as the air-shower energy increases, the stations near the air-shower core have a higher probability of saturation, as mentioned within the context of the water-Cherenkov detectors of the Auger Observatory.
Neither the effects of Poissonian statistics or detector saturation are studied here explicitly, but rather are mentioned as additional considerations further motivating the choice of annulus beyond what is shown in the mass sensitivity curves. 
Therefore, the \qtyrange{800}{850}{\meter} muon annulus was chosen as the nominal annulus for use in this mass sensitivity analysis to compromise between mass sensitivity, and the mentioned qualitative statistical and detector limitations for the studied energy range, spanning several decades in energy.

Additionally,~\Cref{fig:muon_annuli_hadronic_model} shows the proton-iron separation of muon annuli beyond \qty{500}{\meter} from the air-shower axis for QGSJETIII-01 and Sibyll 2.3e simulations.
The resulting FOM curves have a similar relationship with energy as the EPOS LHC-R curves shown in~\Cref{fig:muon_annuli}, highlighting this choice of annulus is independent of hadronic interaction model using the criteria assumed in this analysis.
The QGSJETIII-01 hadronic model generally shows a higher proton-iron mass sensitivity for all muon annuli compared to the other hadronic models, as was also seen and described in~\Cref{chap:simulations}.
As an additional crosscheck for the choice of muon annulus, a past study of simulated \qty[parse-numbers=false]{10^{18.0}}{\electronvolt} vertical air showers ($\theta=\,$\qtyrange{0}{55}{\degree}) showed the optimal distance from the air-shower axis for proton-iron separation of a muon observable is between \qtyrange{600}{800}{\meter}~\cite{holt_epjc2019_radiomuon}.
No detector effects were included in this previous analysis; however, a muon energy threshold of $E{\mu}\,>\,$\qty{1}{\giga\electronvolt} was assumed to obtain the quoted results.
A more detailed study of the muon number mass sensitivity at the location of Auger should be performed at ultra-high-energies ($>$\,\qty[parse-numbers=false]{10^{18.0}}{\electronvolt}) using the muon annulus \qty{1000}{\meter} from the air-shower axis, as this annulus is located exactly at the reference distance used in the Auger SD-1500 reconstruction.

\begin{figure}
  \centering
  \includegraphics[trim={0cm 0cm 0cm 0cm},clip,width=0.45\textwidth]{appendix/B/QGSJETIII01_MassFOM_MuonRings_zen40_50_last10rings_wDataCuts_pFe_electronAtXmaxScaling.pdf}
  \includegraphics[trim={0cm 0cm 0cm 0cm},clip,width=0.45\textwidth]{appendix/B/Sibyll23e_MassFOM_MuonRings_zen40_50_last10rings_wDataCuts_pFe_electronAtXmaxScaling.pdf}
  \caption[Muon annulus merit factors for different hadronic models.]{The proton-iron figure of merit for the muon annuli $>\,$\qty{500}{\meter} from the air-shower axis, studied within a zenith bin of \qtyrange{40}{50}{\degree} for the QGSJETIII-01 (left) and Sibyll 2.3e (right) hadronic interaction models.}
  \label{fig:muon_annuli_hadronic_model}
\end{figure}

%The procedure and analysis of~\Cref{chap:nn} should then be applied to the SD-1500 array, as the results of this thesis show promise for

%The analysis of~\Cref{chap:nn} shows a neural network based reconstruction of station-level muon signals performs reasonably well in modeling the muonic LDF of the WCD measured signals.
%Specifically, the muonic signal at the SD-750 reference distance, as determined from fitting the muon LDF to the neural network predictions, is mass sensitive where the mass sensitivity could increase for larger reference distances. 
%This procedure should then be applied to the SD-1500

%The fit procedure and 

%and therefore, can describe a muonic signal component at a reference distance 

%This muon annulus is nearly equal to or the most mass sensitive muon annulus in most  
%as this annulus is related to $S_{1000}$, the total signal in the Auger WCD \qty{1000}{\meter} from the air-shower axis.

%SD-1500 energy estimator, $S_{1000}$ measured in VEM units. 

%In particular, this is not true in the highest energy bins where the annulus \qtyrange{500}{550}{\meter} from the shower axis provides a larger separation than the \qtyrange{950}{1000}{\meter} annulus

%This analysis spans several decades in air-shower energy, hence a compromise was made in selecting the \qtyrange{800}{850}{\meter} muon annulus as the nominal annulus 

%This effect becomes larger as air shower energy decreases due to their smaller particle footprints on ground, hence

%is larger for lower energy air showers due to their smaller 
%Furthermore, the SD-433 and SD-750 arrays of the Pierre Auger Observatory measure lower energy air showers than the SD-1500 array, with full efficiency for the SD-433 above energies of approximately \qty[parse-numbers=false]{10^{16.8}}{\electronvolt}.
%At these energies, the nominal distance for energy reconstruction 

%the SD-1500 array of the Pierre Auger Observatory uses \qty{1000}{\meter} from the air-shower axis 

%using the FOM metric in~\Cref{eq:fom} to maximize proton-iron separation power while attempting to lessen the ratio of $N_{\mu}$ uncertainty-to-counts. 
%Due to the distribution of low-energy muons at ground and the Poissonian statistics of surface particle detectors then annuli closer to the shower axis will have smaller relative uncertainties.

%The muon number within the farthest annulus exhibits maximum proton-iron separation power, although there are only small differences in separation power between this annulus and other annuli far from the shower axis.
%Furthermore, the SD-1500 Auger array is only fully efficient above 
